\section{\texorpdfstring{\code{volatile}}{volatile}}
\label{sec:volatile}

Hittills har koden bara rört vanligt minne. Den här filen rör \textbf{hårdvara}, och då slutar ett
av kompilatorns grundläggande antaganden att gälla.

Kompilatorn antar att ingen annan än programmet ändrar en variabel. Med den utgångspunkten är
följande loop oändlig, och optimeraren får ta bort den:

\begin{cppcode}
while ((SPI0.INTFLAGS & SPI_IF_bm) == 0U)
{
}
\end{cppcode}

\code{INTFLAGS} läses en gång, resultatet läggs i ett register, och eftersom ingenting i loopen
ändrar det är villkoret konstant. Med \code{-O2} blir resultatet antingen en loop som aldrig
avslutas eller en loop som optimeras bort helt.

Men \code{INTFLAGS} \emph{ändras} --- av periferin, utan att programmet gör något.
\code{volatile} är hur man säger det: \emph{läs om den här variabeln varje gång; slå inte ihop
läsningar, flytta dem inte, ta inte bort dem.}

I \code{<avr/io.h>} är periferiregistren redan deklarerade \code{volatile}. Du behöver inte skriva
kvalificeraren --- men du behöver veta att den är där, och varför.

\subsection{Vad \texorpdfstring{\code{volatile}}{volatile} inte är}

Det är \textbf{inte} en synkroniseringsmekanism. Det garanterar inte atomicitet och ingen ordning
mot annat än andra \code{volatile}-åtkomster. I trådad kod på en värddator är \code{volatile}
nästan alltid fel verktyg; \code{std::atomic} är rätt.

Här handlar det om något annat: ett register som en periferi ändrar under näsan på programmet. Det
är precis det fall \code{volatile} finns för.

\subsection{Var det får stå}

\textbf{Bara i transportens implementationsfil.}

Det är hela poängen med sömmen i \kapref{ch:byte}. Allt ovanför --- \code{Spi}, \code{Stub},
\code{Interface}, \code{EchoNode} --- är vanlig, portabel C++ som kompilerar och körs på din
laptop. En granskare som ser \code{volatile} i någon annan fil ska fråga varför.

\section{SPI-periferin}
\label{sec:periph}

Fyra register räcker för en polld master: ett för enable och klockinställningar, ett för läge, ett
för statusflaggor, och ett för data.

\begin{cppcode}
void AvrSpiTransport::init() noexcept
{
    // Route SPI0 to PORTC (ALT1): PC0 = MOSI, PC1 = MISO, PC2 = SCK, PC3 = SS.
    PORTMUX.SPIROUTEA = PORTMUX_SPI0_ALT1_gc;

    // MOSI, SCK and SS are outputs; MISO is an input. SS idles high.
    PORTC.DIRSET = PIN0_bm | PIN2_bm | PIN3_bm;
    PORTC.DIRCLR = PIN1_bm;
    PORTC.OUTSET = PIN3_bm;

    // Master, MSB first, 4 MHz / 4 = 1 MHz, enabled.
    SPI0.CTRLA = SPI_MASTER_bm | SPI_PRESC_DIV4_gc | SPI_ENABLE_bm;

    // SPI mode 0, and software-controlled slave select.
    SPI0.CTRLB = SPI_MODE_0_gc | SPI_SSD_bm;
}
\end{cppcode}

Fem rader, och tre av dem är fällor:

\textbf{Pinnmuxen.} Utan den raden ligger periferin på sin standardmux, och de fyra ledningar du
kopplat är helt tysta. Symptomet är det mest förvirrande som finns: programmet kör, inga fel
rapporteras, och logikanalysatorn visar ingenting alls. \textbf{Sätt den först.}

\textbf{Slave select disable.} Utan den övervakar periferin \code{SS}-pinnen, och en låg nivå där
får mastern att tro att en annan master tagit över bussen: den lämnar masterläget. Men här är det
\emph{din egen kod} som drar \code{SS} låg vid varje \code{select()}. Utan biten skjuter
drivrutinen sig själv i foten vid första transaktionen.

\textbf{Prescalern.} Protokollet garanterar ingenting över 1~MHz. Kör inte snabbare för att det
går.



\subsection{Transporten}

\begin{cppcode}
void AvrSpiTransport::select() noexcept   { PORTC.OUTCLR = PIN3_bm; }
void AvrSpiTransport::deselect() noexcept { PORTC.OUTSET = PIN3_bm; }

std::uint8_t AvrSpiTransport::transfer(const std::uint8_t byte) noexcept
{
    SPI0.DATA = byte;

    while ((SPI0.INTFLAGS & SPI_IF_bm) == 0U)
    {
    }

    return SPI0.DATA;
}
\end{cppcode}

\textbf{Läsningen av \code{SPI0.DATA} nollställer flaggan.} Hoppa över den och nästa
\code{transfer()} ser en flagga som redan står kvar, returnerar omedelbart, och läser en byte som
inte hunnit komma in. Symptomet är data som ligger \textbf{en byte efter} genom hela
transaktionen, vilket är svårt att se och lätt att missta för ett bitordningsfel.

\textbf{Ingen fördröjning behövs kring \code{SS}.} Protokollet kräver 60~ns setup och hold; en
enda portskrivning på en 4~MHz-MCU är redan flera gånger långsammare än så.

\section{MVIO: två spänningar i samma chip}
\label{sec:mvio}

Det här är den elegantaste detaljen i hela konstruktionen.

\textbf{Problemet:} MCU-kortet körs på 5~V. FPGA-kortets GPIO är 3,3~V och tål inte 5~V. Den
vanliga lösningen är en nivåomvandlare på varje ledning --- en komponent till, fyra kopplingar
till, och fyra saker till som kan sitta fel.

\textbf{Lösningen:} AVR-DB-familjen har \textbf{MVIO} (\emph{Multi-Voltage I/O}). \code{PORTC}
matas inte från \code{VDD} som resten av chippet, utan från en egen pinne: \code{VDDIO2}.

Koppla \code{VDDIO2} till FPGA-kortets 3,3~V-rail, och hela \code{PORTC} blir en
3,3-voltsdomän --- medan kärnan och övriga portar kör oförändrat på 5~V. \textbf{Ingen
nivåomvandlare.}

Det är därför pinnkonfigurationen är en del av kontraktet och inte varje grupps eget val: byter du
till standardmuxen hamnar ledningarna i \code{VDD}-domänen, och då behövs nivåomvandlaren igen
--- eller så går FPGA-ingången sönder.

\begin{aside}
\note \code{VDDIO2} måste vara matad \textbf{innan} \code{PORTC} gör någonting alls. Är den inte
det är portens beteende odefinierat. Statusregistret för MVIO säger om den är det, och en
kontroll av det vid uppstart är den enskilt mest värdefulla raden kod inför bring-upen: en
bortglömd \code{VDDIO2}-koppling ser ut precis som en trasig SPI-implementation, och den
kontrollen skiljer de två fallen åt på en sekund.
\end{aside}

\section{Friställd C++}
\label{sec:freestanding}

Firmware byggs för ett \textbf{friställt} mål. En del av C++ följer inte med:

\begin{booktable}{@{}lL@{}}
\thead{Finns inte} & \thead{Vad du gör i stället} \\ \midrule
Heap (\code{new}, \code{std::vector}) & Statiskt allokerade objekt; allt har känd storlek. \\
Undantag & Returkoder --- vilket hela interfacet redan bygger på. \\
RTTI (\code{dynamic_cast}) & Virtuella funktioner räcker; du castar aldrig nedåt. \\
\code{<iostream>} & UART-utskrift, eller en lysdiod. \\
\end{booktable}

\textbf{Lägg märke till att ingenting av det påverkar designen.} Interfacet returnerar \code{bool}
i stället för att kasta. \code{Frame} är en aggregate med statisk storlek. Ingen klass använder
heapen. Det var inte tur --- projektets kodkonventioner är skrivna för ett friställt mål från
början, och det här är kapitlet där det lönar sig.
